\documentclass[11pt]{article}

\usepackage[margin=1in]{geometry}
\usepackage[T1]{fontenc}
\usepackage[utf8]{inputenc}
\usepackage{lmodern}
\usepackage{microtype}
\usepackage{amsmath,amssymb,amsthm,mathtools}
\usepackage{booktabs,longtable,array}
\usepackage{enumitem}
\usepackage{xcolor}
\usepackage{listings}
\usepackage{hyperref}
\usepackage[nameinlink,noabbrev]{cleveref}

\hypersetup{
  colorlinks=true,
  linkcolor=blue!55!black,
  citecolor=blue!55!black,
  urlcolor=blue!60!black,
  pdftitle={The Goldbach--Euler Theorem in Intuitionistic Primitive Recursive Arithmetic},
  pdfauthor={Research report prepared for Vladimir Reshetnikov}
}

\lstset{
  basicstyle=\ttfamily\small,
  columns=fullflexible,
  frame=single,
  breaklines=true,
  backgroundcolor=\color{black!2},
  keywordstyle=\color{blue!55!black},
  commentstyle=\color{green!40!black}
}

\newtheorem{theorem}{Theorem}[section]
\newtheorem{lemma}[theorem]{Lemma}
\newtheorem{proposition}[theorem]{Proposition}
\newtheorem{corollary}[theorem]{Corollary}
\theoremstyle{definition}
\newtheorem{definition}[theorem]{Definition}
\newtheorem{remark}[theorem]{Remark}
\newtheorem{construction}[theorem]{Construction}

\newcommand{\N}{\mathbb N}
\newcommand{\Q}{\mathbb Q}
\newcommand{\PP}{\mathsf{PP}}
\newcommand{\RF}{\mathsf{RF}}
\newcommand{\Can}{\mathsf{Can}}
\newcommand{\GE}{\mathsf{GE}}
\newcommand{\one}{\mathbf 1}
\newcommand{\set}[1]{\left\{#1\right\}}
\newcommand{\abs}[1]{\left|#1\right|}
\newcommand{\floor}[1]{\left\lfloor#1\right\rfloor}

\title{\textbf{The Goldbach--Euler Theorem in Intuitionistic Primitive Recursive Arithmetic}\\[4pt]
\large A finite rational proof with an explicit primitive-recursive modulus and a Lean formalization roadmap}
\author{Research report prepared for Vladimir Reshetnikov}
\date{15 November 2026}

\begin{document}
\maketitle

\begin{abstract}
The Goldbach--Euler theorem says that the reciprocals of one less than the
\emph{distinct} perfect powers sum to one:
\[
  \frac13+\frac17+\frac18+\frac1{15}+\frac1{24}+\frac1{26}+\cdots=1.
\]
The usual short proof expands each summand as a geometric series, changes the
order of two infinite sums, and telescopes.  Those analytic operations are not
appropriate as primitive steps inside intuitionistic primitive recursive
arithmetic (iPRA).  This report replaces them by a wholly finite argument over
nonnegative rational numbers.  If
\[
 S_N:=\sum_{\substack{2\le p\le N\\ p\text{ a perfect power}}}\frac1{p-1},
\]
where every value $p$ is counted once, then iPRA proves the effective estimate
\[
  0\le 1-S_N\le 2^{-k}
  \qquad\text{whenever}\qquad
  N\ge B(k):=2^{2(k+1)^2}.
\]
Thus the series converges constructively to one with an explicit
primitive-recursive modulus.  The proof uses only bounded perfect-power tests,
finite prime factorization, finite geometric identities, finite sums and
products, and quantifier-free induction.  A staged Lean plan is given for both
a Mathlib-level proof and an internal derivation in the current
\texttt{ConciseIPRA} development.  The archive also includes exact finite
Python checks; these are reproducibility aids and are not presented as formal
verification.
\end{abstract}

\tableofcontents

\section{The theorem and its constructive reading}

The historical theorem is commonly written
\begin{equation}\label{eq:classical-ge}
  \sum_{p}\frac1{p-1}=1,
\end{equation}
where $p$ ranges over perfect powers larger than one, with repetitions omitted.
Thus $16=2^4=4^2$ contributes the single term $1/15$.  Euler published the
identity in 1737 and attributed it to Goldbach; the historical proof manipulated
the divergent harmonic series.  Modern expositions replace that manipulation
by an absolutely convergent geometric-series argument
\cite{viader2006,zeilberger2006}.

For weak constructive arithmetic, the right statement is not an assertion
about a completed set of real numbers.  It is a uniform family of rational
inequalities.

\begin{definition}[Perfect power]\label{def:pp}
For $n\in\N$, let $\PP(n)$ mean
\[
  \exists a,e\le n\;
  \bigl(2\le a\wedge 2\le e\wedge a^e=n\bigr).
\]
This is a bounded, hence decidable, predicate.  Define
\[
  S_N:=\sum_{n=2}^{N}\frac{\one_{\PP(n)}}{n-1}.
\]
The sum is indexed by the value $n$, so duplicate presentations are omitted by
definition.
\end{definition}

\begin{theorem}[Effective Goldbach--Euler theorem]\label{thm:effective}
Let
\[
  B(k):=2^{2(k+1)^2}.
\]
For all $k,N\in\N$,
\[
  B(k)\le N
  \quad\Longrightarrow\quad
  0\le 1-S_N\le 2^{-k}.
\]
Every function and relation in this assertion is primitive recursive, and the
assertion is provable in iPRA.
\end{theorem}

The ordinary theorem follows immediately: the sequence $(S_N)$ has limit one,
with $B$ as a dyadic modulus.  More importantly, \cref{thm:effective} itself is
the internal constructive content.  It is stronger than a bare convergence
claim because it computes a sufficient cutoff from the requested accuracy.

\subsection{A quotient-free arithmetical formulation}

No quotient type for rational numbers is required.  Define primitive-recursive
natural numbers $A(N),D(N)$ by accumulating the finite sum.  Start with
$A(1)=0$, $D(1)=1$.  On passing from $n$ to $n+1$, leave the pair unchanged if
$\neg\PP(n+1)$; if $\PP(n+1)$, put
\[
  A(n+1)=nA(n)+D(n),\qquad D(n+1)=nD(n).
\]
Then $D(N)>0$ and $S_N=A(N)/D(N)$.  The conclusion of
\cref{thm:effective} is equivalent to the natural-number formula
\begin{equation}\label{eq:natural-formula}
  A(N)\le D(N)
  \quad\wedge\quad
  (2^k-1)D(N)\le 2^k A(N).
\end{equation}
Both inequalities are expressed using addition and multiplication only.  This
is the version to be represented by a closed formula of the ordinary
\texttt{PRAAx} language in the Lean development.

\section{Why the short analytic proof must be finitized}

Let $R$ be the integers $r\ge2$ that are not themselves perfect powers.  The
standard short proof observes that every integer $m\ge2$ has a unique
representation $m=r^i$ with $r\in R$ and $i\ge1$, and computes
\[
\begin{aligned}
 \sum_{r\in R}\sum_{s\ge2}\frac1{r^s-1}
 &=\sum_{r\in R}\sum_{s\ge2}\sum_{i\ge1}r^{-si}\\
 &=\sum_{r\in R}\sum_{i\ge1}\sum_{s\ge2}(r^i)^{-s}\\
 &=\sum_{m\ge2}\frac1{m(m-1)}=1.
\end{aligned}
\]
This is mathematically sound in ordinary analysis because every term is
nonnegative.  It is nevertheless an inconvenient proof object for iPRA: it
mentions completed infinite sums, invokes an infinite reindexing, and silently
uses convergence theorems about rational or real sequences.

The argument below keeps the same combinatorial insight but uses it only on
finite rectangles.  The upper bound is proved by taking a sufficiently long
\emph{finite} prefix of each geometric expansion and deriving a contradiction
from a finite telescoping sum.  The lower bound embeds a finite telescoping
rectangle into $S_{B(k)}$.  No limit theorem is used inside the arithmetic.

\section{Canonical roots of integers}

\begin{definition}[Root-free base]\label{def:rootfree}
For $r\ge2$, write $\RF(r)$ if $r$ is not a perfect power.  A
\emph{canonical representation} of $n\ge2$ is an equality
\[
  n=r^e
\]
with $\RF(r)$ and $e\ge1$.
\end{definition}

\begin{lemma}[Canonical root decomposition]\label{lem:canonical}
Every $n\ge2$ has exactly one canonical representation $n=r^e$.
Moreover, $r$ and $e$ are primitive-recursive functions of $n$.
\end{lemma}

\begin{proof}
Write the finite prime factorization
\[
  n=\prod_{j<t}q_j^{\alpha_j},
\]
where the $q_j$ are distinct primes and every $\alpha_j>0$.  Let
\[
  e:=\gcd(\alpha_0,\ldots,\alpha_{t-1}),
  \qquad
  r:=\prod_{j<t}q_j^{\alpha_j/e}.
\]
Then $n=r^e$.  The exponents in the prime factorization of $r$ have greatest
common divisor one.  Hence $r$ cannot equal $u^d$ with $d\ge2$, because then
every prime exponent of $r$ would be divisible by $d$.  Thus $\RF(r)$.

For uniqueness, suppose also $n=u^d$ with $\RF(u)$.  Write
$u=\prod q_j^{\beta_j}$, inserting zero exponents if necessary.  Unique
factorization gives $\alpha_j=d\beta_j$ for every $j$.  Since $u$ is root-free,
the positive $\beta_j$ have gcd one; otherwise their common divisor would
exhibit $u$ as a nontrivial power.  Consequently
\[
  d=\gcd_j\alpha_j=e,
\]
and then $\beta_j=\alpha_j/e$, so $u=r$.

Everything in the construction is bounded.  Trial division computes the prime
factorization; bounded valuation, finite gcd and finite product are primitive
recursive.  Their correctness is proved by quantifier-free induction.  Thus
this is not a use of nonconstructive factor selection.
\end{proof}

\begin{corollary}[Injectivity of canonical powering]\label{cor:injective}
If $\RF(r)$, $\RF(u)$, $i,j\ge1$, and $r^i=u^j$, then $r=u$ and $i=j$.
\end{corollary}

\begin{proof}
Both pairs are canonical representations of their common value, so apply
\cref{lem:canonical}.
\end{proof}

\begin{corollary}[Canonical indexing of perfect powers]\label{cor:perfect-bij}
The map $(r,s)\mapsto r^s$ is a bijection from pairs satisfying
$\RF(r)$ and $s\ge2$ to distinct perfect-power values.
\end{corollary}

\section{Three finite rational identities}

All divisions below are merely notation for nonnegative rational pairs.  Each
identity can be verified after multiplying by a positive common denominator.

\begin{lemma}[Finite denominator expansion]\label{lem:geom-den}
For $r\ge2$, $s\ge1$ and $I\ge1$,
\begin{equation}\label{eq:geom-den}
  \frac1{r^s-1}
  =\sum_{i=1}^{I}\frac1{r^{si}}
   +\frac1{r^{sI}(r^s-1)}.
\end{equation}
In particular, when $s\ge2$ the remainder is at most $2^{-2I}$.
\end{lemma}

\begin{proof}
Put $x=r^s$.  Multiplying the proposed identity by $x^I(x-1)$ reduces it to
\[
  x^I=(x-1)(x^{I-1}+\cdots+1)+1,
\]
the finite geometric identity.  If $r,s\ge2$, then $r^{sI}\ge2^{2I}$ and
$r^s-1\ge1$.
\end{proof}

\begin{lemma}[Finite exponent expansion]\label{lem:geom-exp}
For $m\ge2$ and $T\ge2$,
\begin{equation}\label{eq:geom-exp}
  \sum_{s=2}^{T}\frac1{m^s}
  =\frac1{m(m-1)}-\frac1{m^T(m-1)}.
\end{equation}
\end{lemma}

\begin{proof}
Multiply by $m^T(m-1)$ and telescope, or apply the finite geometric identity
to $m^{-1}$ and then clear denominators.
\end{proof}

\begin{lemma}[Finite telescoping]\label{lem:telescoping}
For $M\ge2$,
\begin{equation}\label{eq:telescoping}
  \sum_{m=2}^{M}\frac1{m(m-1)}=1-\frac1M.
\end{equation}
\end{lemma}

\begin{proof}
Use $1/[m(m-1)]=1/(m-1)-1/m$ and cancel adjacent terms in the finite sum.
\end{proof}

These lemmas are elementary enough to be proved in iPRA by induction on the
finite upper bound.  No appeal to an infinite geometric series has occurred.

\section{The uniform upper bound \texorpdfstring{$S_N\le1$}{S(N) <= 1}}

This is the part of the proof where a careless presentation most easily
reintroduces an infinite sum.  We therefore give the finite argument in full.

For fixed $N$, define the finite canonical index set
\[
  C_N:=\set{(r,s):\RF(r),\ s\ge2,\ r^s\le N}.
\]
By \cref{cor:perfect-bij},
\begin{equation}\label{eq:SN-canonical}
  S_N=\sum_{(r,s)\in C_N}\frac1{r^s-1}.
\end{equation}
Also $|C_N|\le N$.

\begin{proposition}[Finite upper bound]\label{prop:upper}
For every $N$, one has $S_N\le1$.
\end{proposition}

\begin{proof}
The claim is immediate when $C_N$ is empty.  Otherwise choose a positive common
denominator $Q$ for the finite sum in \eqref{eq:SN-canonical}; for example,
\[
  Q:=\prod_{(r,s)\in C_N}(r^s-1).
\]
Write $S_N=P/Q$ with $P\in\N$.  Rational comparison is decidable.  We split on
whether $P\le Q$.  The first case is the desired conclusion.  Suppose instead
that $P>Q$.  Then
\begin{equation}\label{eq:positive-gap}
  S_N\ge1+\frac1Q.
\end{equation}

Set
\[
  I:=NQ+1
\]
and truncate every denominator expansion after $I$ terms:
\[
  A_{N,I}:=
  \sum_{(r,s)\in C_N}\sum_{i=1}^{I}\frac1{r^{si}}.
\]
By \cref{lem:geom-den},
\[
\begin{aligned}
  0\le S_N-A_{N,I}
   &=\sum_{(r,s)\in C_N}
     \frac1{r^{sI}(r^s-1)}\\
   &\le \frac{|C_N|}{2^{2I}}
    \le \frac{N}{2^{2I}}.
\end{aligned}
\]
For every positive integer $x$, $x<2^{2(x+1)}$; this follows by elementary
induction.  With $x=NQ$, we obtain
\[
  \frac{N}{2^{2I}}<\frac1Q.
\]
Combining this with \eqref{eq:positive-gap} gives
\begin{equation}\label{eq:trunc-above-one}
  A_{N,I}>1.
\end{equation}

We now derive the opposite inequality using only a finite regrouping.  For a
fixed root-free $r$ occurring in $C_N$ and a fixed $1\le i\le I$, put
$m=r^i$.  The contribution associated with this pair $(r,i)$ is
\[
  \sum_{\substack{s:(r,s)\in C_N}}\frac1{m^s}.
\]
Every such exponent satisfies $2\le s\le N$: indeed,
$s\le2^s\le r^s\le N$.  Therefore, by the finite geometric identity,
\begin{equation}\label{eq:group-weight}
  \sum_{\substack{s:(r,s)\in C_N}}\frac1{m^s}
  \le\sum_{s=2}^{N}\frac1{m^s}
  <\frac1{m(m-1)}.
\end{equation}
The map $(r,i)\mapsto r^i$ is injective by
\cref{cor:injective}.  Moreover $r\le N$, so $m\le N^I$.  Summing
\eqref{eq:group-weight} over the finitely many pairs $(r,i)$ gives
\[
  A_{N,I}
  <\sum_{m=2}^{N^I}\frac1{m(m-1)}
  =1-\frac1{N^I}<1,
\]
where \cref{lem:telescoping} was used.  This contradicts
\eqref{eq:trunc-above-one}.  Hence the decidable alternative $P>Q$ is
impossible, and $P\le Q$.
\end{proof}

\begin{remark}[Why this reductio is intuitionistically harmless]
The proof does not infer an arbitrary proposition from its double negation.
It computes natural numbers $P,Q$ and performs a decidable comparison.  In the
branch $P>Q$ it derives a contradiction.  Therefore the other branch $P\le Q$
holds.  This is ordinary constructive case analysis on a decidable relation.
\end{remark}

\section{The lower bound and the explicit modulus}

The lower estimate is simpler and exhibits the constructive content directly.
Fix an accuracy parameter $k$ and define
\begin{equation}\label{eq:R-T-B}
  R:=2^{k+1},
  \qquad
  T:=2k+2,
  \qquad
  B:=R^T=2^{2(k+1)^2}.
\end{equation}

For every root-free $r\le R$, let $L_r$ be the largest $i\ge1$ such that
$r^i\le R$.  This is a bounded maximum.  For each $2\le s\le T$,
\cref{lem:geom-den} gives
\[
  \frac1{r^s-1}\ge\sum_{i=1}^{L_r}\frac1{r^{si}}.
\]
Also $r^s\le R^T=B$; by canonical-power injectivity these values are
pairwise distinct, so every term on the left occurs once in $S_B$.  Summing
and then reindexing the finite set by $m=r^i$ yields
\begin{equation}\label{eq:rectangle-embed}
\begin{aligned}
  S_B
  &\ge
  \sum_{\substack{2\le r\le R\\\RF(r)}}
  \sum_{s=2}^{T}\frac1{r^s-1}\\
  &\ge
  \sum_{\substack{2\le r\le R\\\RF(r)}}
  \sum_{s=2}^{T}\sum_{i=1}^{L_r}\frac1{r^{si}}\\
  &=\sum_{m=2}^{R}\sum_{s=2}^{T}\frac1{m^s}.
\end{aligned}
\end{equation}
The last equality is exactly the canonical decomposition lemma: every
$2\le m\le R$ has one and only one form $m=r^i$ with $\RF(r)$.

Now apply \cref{lem:geom-exp,lem:telescoping}:
\begin{align}
 \sum_{m=2}^{R}\sum_{s=2}^{T}\frac1{m^s}
 &=\sum_{m=2}^{R}
   \left(\frac1{m(m-1)}-\frac1{m^T(m-1)}\right)\notag\\
 &=1-\frac1R-E(R,T),\label{eq:lower-exact}
\end{align}
where
\[
  E(R,T):=\sum_{m=2}^{R}\frac1{m^T(m-1)}.
\]
Since $m\ge2$ and $m-1\ge1$,
\begin{equation}\label{eq:E-bound}
  E(R,T)\le\frac{R-1}{2^T}\le\frac{R}{2^T}.
\end{equation}
The choices \eqref{eq:R-T-B} make the two errors equal:
\[
  \frac1R=2^{-k-1},
  \qquad
  \frac{R}{2^T}=2^{k+1-(2k+2)}=2^{-k-1}.
\]
Thus \eqref{eq:rectangle-embed}--\eqref{eq:E-bound} imply
\begin{equation}\label{eq:lower-final}
  S_{B(k)}\ge1-2^{-k}.
\end{equation}

\begin{proof}[Proof of \cref{thm:effective}]
Every summand in $S_N$ is nonnegative, so $S_N$ is monotone in $N$.
By \cref{prop:upper}, $S_N\le1$.  If $N\ge B(k)$, monotonicity and
\eqref{eq:lower-final} give
\[
  1-2^{-k}\le S_{B(k)}\le S_N\le1.
\]
This is exactly the asserted pair of inequalities.
\end{proof}

\section{Formalizability in iPRA}

We now audit every ingredient against the intended foundational system.  Here
iPRA means the ordinary first-order intuitionistic theory whose function
symbols are all primitive-recursive functions and whose induction scheme is
restricted to quantifier-free formulas.  This is the ordinary presentation
implemented by \texttt{PRAAx} in the current ConciseIPRA Lean development
\cite{proveit-pr}.

\subsection{Primitive-recursive data}

The following functions are primitive recursive.
\begin{enumerate}[label=(\arabic*)]
\item Bounded exponentiation, divisibility and equality tests.
\item The Boolean characteristic function of $\PP(n)$, using bounded loops
      over $2\le a,e\le n$.
\item Bounded sums and products, hence the pair $A(N),D(N)$ representing
      $S_N$.
\item Trial-division factorization, or equivalently the bounded prime
      valuation $v_p(n)$ for $p\le n$.
\item Finite gcd of the nonzero valuations and the corresponding canonical
      root and exponent.
\item The bounded maxima $L_r$, finite canonical index sets, and cardinality
      bounds used in the proof.
\item The modulus $B(k)=2^{2(k+1)^2}$.
\end{enumerate}
All loops have explicit input bounds.  No unbounded minimization is present.

\subsection{Induction strength}

The finite geometric and telescoping identities are proved by induction on a
natural upper bound, with all displayed equalities converted to polynomial
identities after clearing positive denominators.  Correctness of trial division,
valuation, finite gcd and finite products is likewise proved by
quantifier-free induction.  More complicated statements are obtained by
universal closure after the relevant quantifier-free inductions have been
performed.

In particular, the proof does not require induction over arbitrary formulas.
The existential content of factorization is supplied by explicit
primitive-recursive output functions, not by an abstract existence theorem.

\subsection{Rational arithmetic is eliminable}

A nonnegative rational is represented by a pair $(a,b)$ with $b>0$.
Equality and order are cross-multiplication:
\[
  (a,b)=(c,d)\iff ad=bc,
  \qquad
  (a,b)\le(c,d)\iff ad\le cb.
\]
Addition and multiplication are primitive recursive on pairs.  Reduction by a
gcd is optional and is not used.  Consequently every line of the proof expands
to a statement about natural numbers.  Formula \eqref{eq:natural-formula} is
the final quotient-free form.

\subsection{Intuitionistic audit}

The proof uses none of the following principles:
\begin{itemize}
\item law of excluded middle for arbitrary propositions;
\item double-negation elimination for an undecidable statement;
\item countable choice;
\item completeness of the real numbers;
\item monotone convergence or Tonelli/Fubini for infinite series;
\item a permutation or regrouping of a completed infinite sum;
\item Markov's principle.
\end{itemize}
Every case distinction is on a bounded Boolean predicate or a natural-number
comparison.  Every reindexing is a bijection between explicitly finite sets.
The only contradiction argument occurs inside a decidable comparison of two
natural numbers representing a rational partial sum.

\subsection{Relation to the current ConciseIPRA repository}

At the inspected ProveIt revision, \texttt{PR.lean} defines primitive-recursive
function symbols \texttt{PR}, terms \texttt{PTm}, and the ordinary iPRA axiom
family \texttt{PRAAx}, including defining equations and quantifier-free
induction.  The file \texttt{NativeBiInterpretation.lean} proves
\texttt{pra\_iff\_native}: an ordinary iPRA derivation is equivalent to a
derivation of its native reading in the scalar Eq/Q/D calculus
\cite{proveit-native}.

Therefore the internal Lean endpoint should have the form
\begin{lstlisting}[language={}] 
theorem goldbachEuler_iPRA :
  ([] : List (Fm PTm 0)) \vdash[PRAAx] goldbachEulerFormula
\end{lstlisting}
followed by
\begin{lstlisting}[language={}] 
(pra_iff_native (Gamma := []) (phi := goldbachEulerFormula)).mp
  goldbachEuler_iPRA
\end{lstlisting}
(up to the repository's Unicode notation).  No new bi-interpretation proof is
needed.

\section{Lean formalization roadmap}

The archive contains two Lean-oriented files and a more detailed roadmap in
\texttt{FORMALIZATION.md}.  They deliberately distinguish a proposed interface
from completed verification.

\subsection{Layer 1: Mathlib theorem}

The first layer should define the bounded perfect-power test and a rational
partial sum indexed by values.  The principal theorem statement is
schematically
\begin{lstlisting}[language={}] 
def modulus (k : Nat) : Nat := 2 ^ (2 * (k + 1)^2)

theorem goldbachEuler_effective (k N : Nat)
    (hN : modulus k <= N) :
    0 <= 1 - partial N /\
    1 - partial N <= (1 : Rat) / (2^k : Nat) := by
  ...
\end{lstlisting}
The proof should be split into the named lemmas already mirrored in
\path{lean/GoldbachEulerOutline.lean}:
canonical decomposition, canonical-power injectivity, two finite geometric
identities, telescoping, the finite upper bound, the modulus lower bound, and
monotonicity.

For the canonical decomposition there are two reasonable implementation
routes.
\begin{enumerate}[label=(\alph*)]
\item Use Mathlib's natural factorization and finitely supported prime
      exponent functions, then prove the gcd characterization of the
      canonical exponent.
\item Implement bounded trial division directly.  This is longer at the
      Mathlib level but aligns almost line by line with the later iPRA coding.
\end{enumerate}
Route (a) is preferable for the semantic theorem; route (b) is preferable for
the internal proof object.  A refinement theorem can connect them.

\subsection{Layer 2: primitive-recursive codes}

Construct \texttt{PR} codes for all arithmetic used in
\eqref{eq:natural-formula}.  It is useful to organize these as a small library:
Boolean choice, bounded folds, exact division under a divisibility premise,
prime valuation, factor-vector gcd, and rational-pair accumulation.

Each code should come with a graph or evaluation theorem.  The graph formulas
remain quantifier-free when output witnesses are supplied by terms.  This is
the form best suited to \texttt{PRA.qfInd}.

\subsection{Layer 3: internal proof}

Prove the natural-number analogue of each finite lemma.  The upper-bound proof
should use the explicit index
\[
  I=N D(N)+1,
\]
so that no existential search for a sufficiently large truncation is needed.
The lower-bound proof uses exactly
\[
  R=2^{k+1},\qquad T=2k+2.
\]
Thus the formal proof can be a direct composition of named estimates rather
than a general development of infinite series.

\subsection{Layer 4: native transfer}

Apply \texttt{pra\_iff\_native}.  The included file
\texttt{lean/ConciseIPRABridge.lean} gives the generic closed-formula transfer
lemma.  This final step is already supported by the current repository; the
substantial remaining work is the arithmetic derivation in \texttt{PRAAx}.

\section{Exact finite validation artifacts}

The accompanying program \texttt{checks/finite\_check.py} uses only the Python
standard library and exact \texttt{Fraction} arithmetic.  It performs the
following audits.
\begin{enumerate}[label=(\arabic*)]
\item It computes the canonical root by prime-exponent gcd and brute-force
      checks uniqueness through $n=3000$.
\item It checks the finite denominator and exponent identities and the
      telescoping formula exactly.
\item For $0\le k\le7$, it checks the finite rectangle inequality
      $T(R,T)\ge U(R,T)\ge1-2^{-k}$.
\item It checks sample upper truncations and direct modulus instances for
      $k=0,1,2$.
\end{enumerate}
A sample from the generated report is shown in \cref{tab:checks}.

\begin{table}[ht]
\centering
\small
\begin{tabular}{rrrrr}
\toprule
$k$ & $R$ & $T$ & finite lower rectangle $U(R,T)$ & $1-2^{-k}$\\
\midrule
0 & 2   & 2  & 0.2500000000 & 0.0000000000\\
1 & 4   & 4  & 0.6800250772 & 0.5000000000\\
2 & 8   & 6  & 0.8585855004 & 0.7500000000\\
3 & 16  & 8  & 0.9335116545 & 0.8750000000\\
4 & 32  & 10 & 0.9677646224 & 0.9375000000\\
5 & 64  & 12 & 0.9841298975 & 0.9687500000\\
6 & 128 & 14 & 0.9921263590 & 0.9843750000\\
7 & 256 & 16 & 0.9960784795 & 0.9921875000\\
\bottomrule
\end{tabular}
\caption{Exact checks are performed with rational arithmetic; decimals are
shown only for readability.}
\label{tab:checks}
\end{table}

These computations are not part of the logical proof.  They detect indexing,
off-by-one and duplicate-counting errors, and they provide reproducible test
vectors for a future Lean implementation.

\section{Comparison with the familiar proof}

The classical one-line argument and the present proof share one decisive idea:
canonical roots convert the two-dimensional family of geometric terms into the
one-dimensional telescoping kernel $1/[m(m-1)]$.  The difference is logical
packaging.

\begin{center}
\begin{tabular}{p{0.42\linewidth}p{0.48\linewidth}}
\toprule
Familiar analytic presentation & Present iPRA presentation\\
\midrule
Infinite geometric expansion & Finite identity with an explicit remainder\\
Interchange of nonnegative sums & Reindexing of an explicitly finite set\\
Infinite telescoping limit & Finite telescoping identity plus a rational error\\
Existence of a sufficiently large cutoff & Closed primitive-recursive formula $B(k)$\\
Real convergence & Natural-number inequalities for numerator/denominator pairs\\
\bottomrule
\end{tabular}
\end{center}

The explicit modulus is not merely an accommodation to weak arithmetic.  It
reveals a quantitative statement hidden by the standard proof.  The bound is
not optimized, but its simple square exponent makes the formal proof short.

\section{Further research questions}

\begin{enumerate}[label=\textbf{Q\arabic*.},leftmargin=*]
\item \textbf{Sharper primitive-recursive modulus.}
Can one retain a comparably short iPRA proof while replacing
$2^{2(k+1)^2}$ by a substantially smaller cutoff?  The proof deliberately
uses very crude bounds $m^T(m-1)\ge2^T$ and $|C_N|\le N$.

\item \textbf{Proof without prime factorization.}
Can canonical roots be developed in iPRA from a direct common-base lemma for
$r^i=u^j$, avoiding an explicit factorization library?  Such a proof might be
shorter in the native iterator language even if it is longer mathematically.

\item \textbf{Direct native proof.}
Is it more economical to prove \eqref{eq:natural-formula} first in
\texttt{PRAAx} and transfer it, or to construct the proof directly in one of
the smaller native clause calculi?  The latter could provide a useful
benchmark for the practical proof complexity of ConciseIPRA.

\item \textbf{Reusable finite-series library.}
The pattern ``finite geometric expansion + explicit remainder + finite
reindexing'' applies to many elementary series identities.  A library of such
lemmas could support constructive formalizations of Euler products, Lambert
series and elementary generating-function arguments without developing real
analysis.

\item \textbf{Certificate extraction.}
Can the eventual Lean proof export a compact ConciseIPRA certificate checked by
one of the repository's small standalone proof checkers?  This would connect a
substantive number-theoretic theorem with the project's proof-checker-size
metric.
\end{enumerate}

\appendix

\section{Detailed primitive-recursive implementation sketch}

This appendix records one concrete coding route.  It is not needed for the
human proof, but it removes ambiguity from the claim of iPRA formalizability.

\subsection{Bounded arithmetic primitives}

Use primitive recursion to define addition, multiplication, exponentiation,
truncated subtraction, sign, comparison, Boolean conjunction and conditional
selection.  For a Boolean predicate $b(i,\vec x)$, define bounded sums and
products by
\[
\begin{aligned}
 \operatorname{sum}(0,\vec x)&=0,
 &\operatorname{sum}(n+1,\vec x)&=\operatorname{sum}(n,\vec x)+f(n,\vec x),\\
 \operatorname{prod}(0,\vec x)&=1,
 &\operatorname{prod}(n+1,\vec x)&=\operatorname{prod}(n,\vec x)f(n,\vec x).
\end{aligned}
\]
Bounded existential and universal tests are folds of Boolean disjunction and
conjunction.

\subsection{Perfect powers}

Define
\[
 \chi_{\PP}(n)=
 \bigvee_{a=2}^{n}\bigvee_{e=2}^{n}[a^e=n].
\]
This is primitive recursive and extensionally equivalent to
\cref{def:pp}.  No logarithm is required; the wasteful bound $e\le n$ is
convenient in weak arithmetic.

\subsection{Prime valuations and canonical roots}

Define primality by bounded trial division.  Define $v_p(n)$ as the largest
$e\le n$ with $p^e\mid n$, using a bounded maximum.  Let
\[
  g(n)=\gcd\set{v_p(n):2\le p\le n,\ p\text{ prime},\ v_p(n)>0}.
\]
The set is nonempty for $n\ge2$.  A fold with the convention
$\gcd(0,a)=a$ computes $g(n)$.  Then define
\[
  \operatorname{root}(n)=
  \prod_{\substack{2\le p\le n\\p\text{ prime}}}
  p^{v_p(n)/g(n)}.
\]
The quotient is exact.  Bounded division is primitive recursive even without
an exactness premise; correctness supplies the premise in the theorem.

An alternative implementation repeatedly divides by the least divisor and
stores exponents in a bounded sequence code.  The valuation version is simpler
to state, while the repeated-division version may generate smaller terms.

\subsection{Rational-pair accumulation}

Use the recurrence preceding \eqref{eq:natural-formula}.  Its invariant is
proved by induction:
\[
  \frac{A(N)}{D(N)}=
  \sum_{n=2}^{N}\frac{\chi_{\PP}(n)}{n-1}.
\]
Internally, state the invariant after clearing all denominators.  Positivity of
$D(N)$ follows simultaneously.

\subsection{Encoding the main theorem}

Let $\operatorname{le}(x,y)$ abbreviate a formula asserting
$\exists z\,x+z=y$, or use a primitive-recursive Boolean comparison and its
graph.  The closed theorem is the universal closure of
\[
  B(k)\le N\to
  \left(A(N)\le D(N)\wedge
  (2^k-1)D(N)\le2^kA(N)\right).
\]
Although this formula has quantifiers and connectives, each induction used in
its proof is on a quantifier-free invariant.  This is exactly the division of
labor supported by ordinary iPRA.

\section{Proof dependency graph}

\begin{center}
\begin{tabular}{p{0.27\linewidth}p{0.65\linewidth}}
\toprule
Node & Depends on / supplies\\
\midrule
Bounded arithmetic & primitive recursion; supplies powers, folds, comparisons\\
Finite factorization & bounded arithmetic; supplies valuations and unique factor data\\
Canonical root lemma & factorization and finite gcd; supplies injective map $(r,i)\mapsto r^i$\\
Finite geometric lemmas & bounded arithmetic; supplies exact prefix/remainder formulas\\
Finite telescoping & bounded arithmetic; supplies $1-1/M$ identity\\
Upper bound & canonical injectivity, denominator expansion, telescoping\\
Lower rectangle & canonical bijection, both geometric identities, telescoping\\
Effective theorem & upper bound, lower rectangle, monotonicity\\
Natural-pair theorem & rational elimination and primitive-recursive coding\\
Native ConciseIPRA theorem & \texttt{pra\_iff\_native}\\
\bottomrule
\end{tabular}
\end{center}

\section{Status ledger}

\begin{itemize}
\item \textbf{Human mathematics:} complete proof of the effective theorem,
      including explicit constants and a constructive audit.
\item \textbf{iPRA reduction:} detailed primitive-recursive encoding and a
      closed quotient-free target formula; no principle beyond iPRA is used.
\item \textbf{Lean:} theorem interfaces and a dependency-complete roadmap are
      supplied.  The Mathlib outline contains \texttt{sorry} and is not
      claimed to compile or to be verified.
\item \textbf{ConciseIPRA bridge:} the final generic transfer uses the existing
      proved theorem \path{pra_iff_native}; construction of the full
      internal arithmetic derivation remains to be implemented.
\item \textbf{Executable checks:} exact finite audits pass, but they are not a
      substitute for formal proof.
\end{itemize}

\begin{thebibliography}{9}

\bibitem{viader2006}
Pelegr\'i Viader, Llu\'is Bibiloni, and Jaume Parad\'is,
\emph{On a Series of Goldbach and Euler},
American Mathematical Monthly \textbf{113} (2006), 206--220.
\href{https://doi.org/10.2307/27641889}{doi:10.2307/27641889}.

\bibitem{zeilberger2006}
Doron Zeilberger,
\emph{A Proof of the Celebrated Goldbach's Theorem}, 2006.
\url{https://sites.math.rutgers.edu/~zeilberg/goldbach.pdf}.

\bibitem{wikipedia-ge}
\emph{Goldbach--Euler theorem}, Wikipedia.
\href{https://en.wikipedia.org/wiki/Goldbach\%E2\%80\%93Euler_theorem}{Wikipedia article}
(accessed 15 November 2026).

\bibitem{proveit-readme}
Vladimir Reshetnikov,
\emph{Concise presentations of intuitionistic primitive recursive arithmetic},
ProveIt repository, revision
\texttt{a85c645dd7213a1adbe35b0965d1e44e06c942be}.
\href{https://github.com/VladimirReshetnikov/ProveIt/tree/main/Logic/PrimitiveRecursiveArithmetic/ConciseIPRA}{ConciseIPRA project directory}.

\bibitem{proveit-pr}
Vladimir Reshetnikov,
\texttt{ConciseIPRA/PR.lean}: primitive-recursive function symbols, terms and
\texttt{PRAAx}, same repository revision.

\bibitem{proveit-native}
Vladimir Reshetnikov,
\texttt{NativeBiInterpretation.lean}: faithfulness,
bi-interpretation and \texttt{pra\_iff\_native}, same repository revision.

\end{thebibliography}

\end{document}
